\sscp{Pronominal sets}{13-03-02}
There is wide variation in the forms,
functions, and distribution of pronominal sets in Austronesian
languages in the region.

The most historically stable forms in the region seem to relate
to first person (singular
and plural) forms: PMP *i-aku ‘\tsc{1sg} free’,
*ku- ‘\tsc{1sg.actor}’, *-ku ‘\tsc{1sg.gen}’,
*kita ‘\tsc{1pl.incl} free’, *ta- ‘\tsc{1pl.incl.actor}’,
*kami ‘\tsc{1pl.excl} free’, *mi- ‘\tsc{1pl.excl.actor}’
\citep{AdelaarHajek2024}.
Even within these retentions there is considerable variation,
along with divergent and unrelated forms in the same subgroups and nodes.

Second and third person (singular
and plural) forms show considerable variation,
such that retentions of forms such as: *i-kahu ‘\tsc{2sg} (free pronoun)’,
*mu- ‘\tsc{2sg.actor}’, *-mu ‘\tsc{2sg.gen}’,
*kimuyu/*kamuyu ‘\tsc{2pl} (free pronoun)’, *ni-a ‘\tsc{3sg} (Actor,
Undergoer, Genitive)’, *sida ‘\tsc{3pl} (free pronoun)’, *-da ‘\tsc{3pl.gen}’
are retained sporadically, but are also interspersed with other divergent
and unrelated forms in the same subgroups and nodes.

Pronominal data has been provided from different languages in
several subgroups in preceding chapters.
Some additional data are given here to illustrate the wide variation
in pronominal systems in the region,
and to show that what is similar in form needs to
be treated as historical retentions,
rather than innovations, and thus are not useful for subgrouping.

\begin{table}\tcp{Comparing Havu and Dhao pronominal systems}{13-14}
%\begin{adjustbox}{width=\textwidth}
	\begin{threeparttable}\st{6pt}
	\begin{tabular}{lllll|ll}\lsptoprule
								&	\mc{4}{l|}{Dhao}															&	Havu													\\
								&	Subject	&	Undergoer	&	Clitic		&	Free	{\cdr}&	Free	{\cdr}&	variants\su{Δ}	\\	\midrule
\tsc{1sg}				&	k-, ku-	&	-ku				&	ku				&	ʤaʔa	{\cdr}&	yaa		{\cdr}&	ʄaa, ʤoo\\
\tsc{2sg}				&	m-, mu-	&	-mu				&	mu				&	əu		{\cdr}&	əu		{\cdr}&	\\
\tsc{3sg}				&	n-			&	-ʔe				&	na\su{†}	&	nəŋu	{\cdr}&	noo		{\cdr}&	\\
\tsc{1pl.incl}	&	t-			&	-ti				&	ti				&	əɖ͡ʐi	{\cdr}&	dii		{\cdr}&	\\
\tsc{1pl.excl}	&	ŋ-			&	-ʔa				&	ŋa				&	ʤiʔi	{\cdr}&	ʄii		{\cdr}&	\\
\tsc{2pl}				&	m-, mi-	&	-mi				&	mi				&	miu		{\cdr}&	muu		{\cdr}&	\\
\tsc{3pl}				&	r-			&	-si				&	ra\su{‡}	&	rəŋu	{\cdr}&	roo		{\cdr}&	noo\su{ǁ}\\
\lspbottomrule
	\end{tabular}
		\begin{tablenotes}\tnsize
			\item[†]	{An alternate \tsc{3sg} object enclitic which adds
								proximal referential deixis is \it{ne}.}
			\item[‡]	{An alternate form for \tsc{3pl} object enclitic which adds
								proximal referential deixis is \it{si}.
								\it{si} further functions as an associative plural.}
			\item[Δ]	{The dominant Seba dialect uses \tsc{1sg} \it{yaa}.
								Dimu and Raijua Wawa use \it{ʄaa}.
								Raijua D{\Q}ida uses \it{ʤoo}. \tsc{3pl} \it{noo} is Raijua.}
		\end{tablenotes}
	\end{threeparttable}
%\end{adjustbox}
\end{table}

Closely related \tsc{Bima-Lembata} languages Havu
and Dhao use their free pronouns as subject,
object, object of preposition, and possessor (see \trf{02-02})
for the latter, using post-posed possessors and showing no alienable-inalienable contrast).
Dhao clitics can be used as proclitics or more commonly as enclitics.
They can optionally be an enclitic to mark post-posed possessor
or object, or some speakers occasionally use them as proclitics for preposed subject.
Many Dhao verbs are not inflected for subject.
A few verbs are prefixed with subject-indexing.
A few of those have irregular inflection reflecting traces of
an underlying historical vowel.
Only one verb of motion ‘go’ is inflected with Undergoer suffixes.
The data in \trf{13-14} are adapted from \citet[264ff]{Grimes2010a}.

The pronominal systems of several \tsc{Ambon-Seram} languages
and Banda were illustrated in \crf{ch:CM},
showing multiple inflectional paradigms each with different morphophonemic
alternation on the initial consonant of the roots (see Tables
\ref{tab:07-16} and \ref{tab:07-17},  several Split-S languages with Actor proclitics
and Undergoer enclitics, and a human/non-human distinction in
third person pronouns (see \trf{07-18}).
By contrast, nearby \tsc{Sula-Buru} languages seem to show none
of these features, including no verbs that are inflected for person and number.
\tsc{Sula-Buru} languages are also very different to the \tsc{Havu-Dhao}
forms, functions, and distribution in \trf{13-14}.

\begin{table}\tcp{Buru pronominal sets}{13-15}
\begin{adjustbox}{width=\textwidth}
	\begin{threeparttable}\st{3pt}
	\begin{tabular}{llllllllll}\lsptoprule
							&	Free	&	Cliticised	&	\mc{2}{l}{Pronominal}	&	\mc{3}{l}{Possessive word}	&	Genitive\su{†}	&	Archaic	\\
							&				&	actor				&	\mc{2}{l}{proclitics}	&	\mc{3}{l}{`have, own, associated'}	&	enclitics	&	undergoer	\\	\midrule
0							&				&							&							&					&	nain								&&		&	-t	&		\\
\tsc{1sg}			&	yako	&	yak, yaʔ, ya&	a						&					&	naŋ(o)							&&		&	-ŋ	&	kon(o)	\\
\tsc{2sg}			&	kae		&	ka					&	ku		{\cdr}&					&	nam(o)				{\cdr}&&		&	-m(o)	&		\\
\tsc{3sg}			&	riŋe	&	riŋ					&	da, d				&					&	\mc{3}{l}{nak(e), naʔ}		&	-n	&		\\
\tsc{1pl.incl}&	kami	&	kam					&	ma		{\cbl}&					&	nam(i)				{\cdr}&&		&	-nam	&		\\
\tsc{1pl.excl}&	kita	&	kit					&	ma		{\cbl}&					&	nan(i)							&&		&	-nan	&		\\
\tsc{2pl}			&	kimi	&	kim					&	ku		{\cdr}&					&	nim(i)							&&		&	-nim	&		\\
\tsc{3pl}			&	sira	&	sir					&	du, d				&					&	\mc{3}{l}{nin(i)/nun(u)\su{‡}}	&	-nin	&		\\
\tsc{3dl}			&	siro	&							&							&					&											&&		&		&		\\
\lspbottomrule
	\end{tabular}
		\begin{tablenotes}\tnsize
			\item[†]	{In the Masarete dialect the genitive system has collapsed to
								the \tsc{3sg} \it{-n} form for all person and number combinations.
								\citet{Collins1980} describes a similar collapse of the genitive
								system for Laha, spoken around the airport on Ambon.}
			\item[‡]	{The Masarete dialect of Buru uses \it{nunu}.
								Rana and Wae Sama use \it{nini}.
								Related Lisela also uses \it{nini}.}
		\end{tablenotes}
	\end{threeparttable}
\end{adjustbox}
\end{table}

The Buru data in \trf{13-15} are adapted from \citet{Grimes1991c}.
Note that the inclusive-exclusive distinction is neutralised
in the \tsc{1pl} pronominal proclitics.
The singular-plural distinction is similarly neutralised in the
pronominal proclitics between \tsc{2sg} and \tsc{2pl}.
The short forms of the possessive word similarly show no contrast
between \tsc{2sg} and \tsc{1pl.excl}.
Free pronouns can be used for subject,
object, object of a preposition, and possessor.
The cliticised pronouns can be used for subject
and possessor (interacting optionally with both the possessive
word and the genitive enclitics in the two possessive constructions;
see \srf{02-03-01-01}), Pronominal proclitics can only be subject,
but semantically can be either Actor or Undergoer for active and non-active verbs.
The Undergoer \it{kono} only ever follows the verb.
It is still productive in the Rana dialect,
but obsolete in the Masarete dialect.

The \tsc{Seram-Tanimbar-Bomberai} language Kowiai similarly neutralises
the singular-plural distinction for second
and third person forms in several categories,
as shaded in \trf{13-16}.
It also makes an alienable-inalienable contrast in the possessive system,
but further marks some inalienable nouns with prefixes,
and others with suffixes.
Data in \trf{13-16} are gleaned from \citet{WalkerWalker1991}.

\begin{table}
	\centering\tcp{Kowiai prominal sets}{13-16}
	\st{6pt}\begin{tabular}{llllll}\lsptoprule
	&	Free	&	Active verb	&	Inalienable	&	Inalienable 	&	Alienable \\
	&	pronouns	&	\tsc{subj} prefixes	&	prefixes	&	suffixes	&	possessive\\	\midrule
\tsc{1sg}				&	laʔ		&	u-	&	nuŋgu-			&	-yoŋ	&	laʔo\\
\tsc{2sg}				&	au		&	mu-	{\cdr}&	mu-					{\cdr}&	-yom	{\cdr}&	aɸo\\
\tsc{3sg}				&	i			&	na-	{\cbl}&	ni-					{\cbl}&	{\0}	{\cbl}&	i-ɸo\\
\tsc{1pl.incl}	&	ʔam		&	a-	&	am nuŋgu-		&	-mam	&	ambo\\
\tsc{1pl.excl}	&	ʔita	&	ta-	&	ita nuŋgu-	&	-nina	&	ita-ɸo\\
\tsc{2pl}				&	o			&	mu-	{\cdr}&	mu-					{\cdr}&	-yom	{\cdr}&	o-ɸo\\
\tsc{3pl}				&	si		&	na-	{\cbl}&	ni-					{\cbl}&	{\0}	{\cbl}&	si-ɸo\\
		\lspbottomrule
	\end{tabular}
\end{table}

\largerpage
In \tsc{Timor-Babar} languages of the \tsc{Rote-Meto} node,
for verbs that do take subject indexing,
a distinction is found between syllabic
and non-syllabic verb agreement prefixes.
In \trf{04-06} we showed that patterned distinctions in the forms
of the syllabic verb agreement prefixes are part of the evidence
justifying a division between the \tsc{Nuclear Rote}
and the \tsc{West Rote-Meto} nodes.

There is a development in the morphosyntax that also sets \tsc{Rote-Meto}
languages apart from surrounding languages in nearby nodes such
as Helong (\srf{04-01}, see also \citealp[75]{Edwards2020}),
\tsc{Eastern Timor} languages such as Tetun (\srf{04-03}),
\tsc{Bima-Lembata} languages such as Havu
and Dhao (\srf{03-03}, see also \citealp[264]{Grimes2010a}),
or \tsc{Central Timor} languages such as Kemak (\crf{ch:CT}).
A contrast has developed,
particularly in the \tsc{West Rote-Meto} pronouns between
accusative and non-accusative pronouns, with the accusative forms also being
used in oblique roles.
There are additional complexities in \tsc{Rote-Meto} pronominal
systems not presented in the next three tables,
since the focus here is limited to the development of a distinction
between nominative (or non-accusative) and accusative pronouns.
In the \tsc{West Rote-Meto} languages,
many non-accusative forms begin with {\#}h,
even if their equivalent historical etymon begins with *k.
But once again, what is cognate in \tsc{Rote-Meto} languages
with languages in other nodes
and subgroups are historical retentions,
not innovations, and therefore we find nothing to link them through
innovations with other Wallacean subgroups.

Pronouns from Meto languages are illustrated in \trf{13-17} by
data from Kotos Amarasi
(adapted from \citealp{GrimesBaniCaet2012} and \citealp[262]{Edwards2020}).

\begin{table}
	\centering\tcp{\tsc{Meto} (Kotos Amarasi) pronominals}{13-17}
	\st{6pt}\begin{tabular}{lll|ll|l}\lsptoprule
&	\tsc{nom/poss}	&	\tsc{acc}	&	\mc{2}{l|}{\tsc{subj} prefix}	&	\tsc{gen} suffix	\\	\midrule
\tsc{1sg}	&	au	&	=kau	&	u-	&	ʔ-	&	-k	\\
\tsc{2sg}	&	hoo	&	=koo	&	mu-	&	m-	&	-m	\\
\tsc{3sg}	&	iin/ini	&	=ee	&	na-	&	n-	&	-n	\\
\tsc{1pl.incl}	&	hit/hiti	&	=kiit/=kiti	&	ta-	&	t-	&	-k	\\
\tsc{1pl.excl}	&	hai	&	=kai	&	mi-	&	m-	&	-m	\\
\tsc{2pl}	&	hii	&	=kii	&	mi-	&	m-	&	-m	\\
\tsc{3pl}	&	siin/sini	&	=sin	&	na-	&	n-	&	-k	\\
\lspbottomrule
	\end{tabular}
\end{table}

In the typology of alignment of S,
A and O \citep{Dixon1994}, Amarasi S
and A both use the nominative pronouns,
as opposed to O,
which uses accusative pronouns.
Because Split-S languages are also common in the wider region,
examples are shown below for non-active intransitive,
active intransitive, and transitive clauses.
As mentioned above, oblique arguments in Amarasi use the accusative
pronouns, as in \qf{13-30}--\qf{13-31}.
The accusative pronouns also function as the object of prepositional
verbs such as \it{nok} ‘with’ or \it{neu} ‘to,
\tsc{dative}’.

\begin{exe}
	\ex{\gll	\tbf{iin} na-meen maʔfenaʔ\\
						\tsc{3sg.nom} 3-sick heavy\\
			\glt	`s/he was seriously ill.'
						}\label{ex:13-27}
	\ex{\gll	\tbf{au} ʔ-nao he ʔ-suub=\tbr{ee}\\
						\tsc{1sg.nom} \tsc{1sg}-go \tsc{irr} \tsc{1sg}-bury=\tsc{3sg.acc}\\
			\glt	`I’m going to bury her/him/it.'
						}\label{ex:13-28}
	\ex{\gll	\tbf{hoo} m-heke m-rair=\tbf{kau}\\
						\tsc{2sg.nom} 2-catch 2-\tsc{pfv=1sg.acc}\\
			\glt	`You have arrested me.'
						}\label{ex:13-29}
	\ex{\gll	iin n-hain=\tbf{kau} afu\\
						\tsc{3sg.nom} 3-dig=\tsc{1sg.acc} soil\\
			\glt	`S/he’s digging soil for me.'
						}\label{ex:13-30}
	\ex{\gll	au ʔ-nao ʔ-roi=\tbf{koo} oe\\
						\tsc{1sg.nom} \tsc{1sg}-go \tsc{1sg}-carry=\tsc{2sg.acc} water\\
			\glt	`I’ll fetch water for you.'
						}\label{ex:13-31}
\end{exe}

In \tsc{Meto},
accusative pronouns can optionally be used in conjunction with
nominative pronouns for certain equative constructions.
Examples are given in \qf{13-32}
and \qf{13-33}.

\begin{exe}
	\ex{\gll	\tbf{au} ka= amna-ah bubur=\tbf{kau} =fa\\
						\tsc{1sg.nom} \tsc{neg1=} \tsc{nmlz}-eat porridge=\tsc{1sg.acc} \tsc{=neg2}\\
			\glt	‘I don’t eat porridge’ literally: ‘I’m not a porridge eater.’
						}\label{ex:13-32}
	\ex{\gll	\tbf{hoo} mun{\tl}munif=\tbf{koo} =t mu-snaas mu-ʔko\\
						\tsc{øsg.nom} \tsc{itns}{\tl}young=\tsc{2sg.acc} =while \tsc{2sg}-stop \tsc{2sg}-from\\
			\glt	‘While you are still young you should stop from (that).’
						}\label{ex:13-33}
\end{exe}

Pronominals from Dela (\tsc{West Rote-Meto}) are illustrated in \trf{13-18}.
Dela does not have a full paradigm of nominative
and accusative pronouns and the caseless pronouns can function
as either subject or object (though the \tsc{2sg},
\tsc{1pl.incl}, and \tsc{2pl} caseless pronouns most often function as subjects).
Dela also has vocative pronouns for \tsc{1pl.incl} and \tsc{2pl}.
The accusative pronouns also function as the object of prepositional
verbs such as \it{noo} ‘with’,
\it{neu} ‘to/for’, or \it{mia} ‘from’.

\begin{table}
	\centering\tcp{Dela pronominals \citep[97]{Tamelan2021}}{13-18}
	\st{5pt}\begin{tabular}{ll|lll|ll|ll}\lsptoprule
								&	free (no case)	&	\tsc{nom}	&	\tsc{acc}	&	\tsc{voc}	&	\mc{2}{l|}{\tsc{subj} prefix}	&	\mc{2}{l}{\tsc{gen} enclitic}\\	\midrule
\tsc{1sg}				&	au	&		&		&		&	ʔu-	&	ʔ-	&	=ŋga	&	=ŋg\\
\tsc{2sg}				&	hoo	&		&	ŋgo	&		&	mu-	&	m-	&	=ma	&	=m\\
\tsc{3sg}				&	eni	&	ana	&	e/ne	&		&	na-	&	n-	&	=na	&	=n\\
\tsc{1pl.incl}	&	hita	&		&	ŋgita	&	ata	&	ta-	&	t-	&	=na/=ta	&	=n\\
\tsc{1pl.excl}	&	hai	&		&		&		&	mi-	&	m-	&	=ma	&	=m\\
\tsc{2pl}				&	hei	&		&	ŋgi	&	ama	&	mi-	&	m-	&	=ma	&	=m\\
\tsc{3pl}				&	sira	&	ara	&	se	&		&	ro-	&	r-	&	=na	&	=n\\
0								&		&		&		&		&	ne-	&		&		&	\\
		\lspbottomrule
	\end{tabular}
\end{table}

In the \tsc{Nuclear Rote} language Rikou,
spoken in eastern Rote,
the distinction between \tsc{nom}
and \tsc{acc} is found only in third person pronouns.
Rikou pronominals are given in \trf{13-19}.
Rikou data are gleaned from \citet{UBB2020b}
and \citet{Jonker1908,Jonker1915}.

\begin{table}
	\centering\tcp{Rikou pronominals}{13-19}
	\st{6pt}\begin{tabular}{ll|ll|ll|l}\lsptoprule
	&	free (no case)	&	\tsc{nom}	&	\tsc{acc}	&	\mc{2}{l|}{\tsc{subj} prefix}	&	\tsc{gen} enclitic\\	\midrule
\tsc{1sg}	&	au	&		&		&	ʔa-	&	ʔ-	&	=ka\\
\tsc{2sg}	&	oo	&		&		&	ma-	&	m-	&	=ma\\
\tsc{3sg}	&	ria	&	ana	&	ni	&	na-	&	n-	&	=na\\
\tsc{1pl.incl}	&	ita	&		&		&	ta-	&	t-	&	=na\\
\tsc{1pl.excl}	&	ami	&		&		&	ma-	&	m-	&	=ma\\
\tsc{2pl}	&	emi	&		&		&	ma-	&	m-	&	=ma\\
\tsc{3pl}	&	sira	&	ara	&	asa, si	&	ra-	&	r-	&	=na\\
		\lspbottomrule
	\end{tabular}
\end{table}

There is variation in Rikou regarding which  form the object
of prepositional verbs surfaces in,
as shown in \trf{13-20},
illustrated with, \it{noo} ‘with’,
\it{fee} ‘to/for’, \it{neuʔ} ‘to/for’,
\it{neme {\ldots} (mai)} ‘from’, and \it{laʔe-neuʔ} ‘about’.
Nominative forms as the object of prepositional verbs are shaded.

\begin{table}
	\centering\tcp{Rikou prepositional verbs}{13-20}
	\st{6pt}\begin{tabular}{llllll}\lsptoprule
&	with	&	from	&	to/for	&	to/for	&	about	\\	\midrule
\tsc{1sg}	&	noo au	&	neme au (mai)	&	fee au	&	neuʔ au	&	laʔe-neuʔ au	\\
\tsc{2sg}	&	noo oo	&	neme oo (mai)	&	fee oo	&	neuʔ oo	&		\\
\tsc{3sg}	&	noo ni	&	neme ana (mai)	{\cdr}&	fee ni	&	neuʔ ana	{\cdr}&	laʔe-neuʔ ana	{\cdr}\\
\tsc{1pl.incl}	&	noo ita	&	neme ita (mai)	&	nfee ita	&	neuʔ ita	&		\\
\tsc{1pl.excl}	&	noo ami	&	neme ami (mai)	&	fee ami	&	neuʔ ami	&	laʔe-neuʔ ami	\\
\tsc{2pl}	&	noo emi	&	neme emi (mai)	&	fee emi	&	neuʔ emi	&	laʔe-neuʔ emi	\\
\tsc{3pl}	&	noo si	&	neme sira (mai)	{\cdr}&	fee si	&	neuʔ asa	&	laʔe-neuʔ asa	\\
\lspbottomrule
	\end{tabular}
\end{table}

As glimpsed in the preceding discussion,
throughout the region there is great diversity in the forms,
functions, and distributions of the pronominal systems,
even within the same subgroup.
Many forms are not cognate.
The pronominal forms that are retentions,
not innovations, have little value for subgrouping purposes.
Subject indexing is discussed in \srf{13-03-03},
noting that some of the forms in different subgroups,
and even within the same subgroup,
are not cognate with each other.
We have not yet been able to identify any innovations in the
pronominal systems that are useful for higher-level subgrouping.
This is also limited by the severe underdocumentation of the
languages in the region.
There is much room for further documentation
and comparative study on this topic.
